% !PRINTED 645
\weberpage{645}
We wish to give one of the radicals, say $\tau_\alpha$, a different sign, that is, to make the permutation

\weq{13}{(\tau_\alpha,\ \varepsilon^\beta \tau_\alpha)}
where $\beta$ can be an arbitrary exponent, and $\alpha$ is one of the indices $0, 1 \ldots n-2$.

In (12) the radical $\tau_\alpha$ has the exponent $r_{n+\nu-\alpha-2} \equiv r_{\nu-\alpha-1}$, and hence to the permutation (13) there corresponds the permutation

\weq{14}{\left(\sqrt[n]{R_\nu},\ \varepsilon^{\beta r_{\nu-\alpha-1}} \sqrt[n]{R_\nu}\right).}
Now, by the meaning of the numbers $r$, we have generally $r_{\alpha+\beta} \equiv r_\alpha r_\beta \pmod{n}$, and if therefore one introduces the permutation (13) into (11), $h$ passes over into $h + \beta r_{-\alpha-1}$.

1. Thus the permutation (13) calls forth among the indices of $\xi$ the cyclic permutation
\[
(h,\ h + \beta r_{-\alpha-1})
\]

If, secondly, we make the cyclic permutation

\weq{15}{(\tau_0,\ \tau_1, \ldots \tau_{n-2}),}
then this corresponds to the permutations

\weq{16}{\left(\sqrt[n]{R_\nu},\ \sqrt[n]{R_{\nu+1}}\right), \quad (K_\nu, K_{\nu+1}), \quad \nu = 0, 1, 2 \ldots n-2,}
and $\xi_h$ passes over into

\[
A + \sum^{\nu} \varepsilon^{h r_\nu} K_{\nu+1} \sqrt[n]{R_{\nu+1}} = A + \sum \varepsilon^{h r_{\nu-1}} K_\nu \sqrt[n]{R_\nu}
\]
\[
= A + \sum \varepsilon^{h g^{-1} r_\nu} K_\nu \sqrt[n]{R_\nu}.
\]

2. Accordingly the indices of $\xi$ suffer, through the permutation (15), the permutation
\[
(h,\ g^{-1} h).
\]

If we now consider any rational symmetric (or even merely metacyclic) function of the $\xi$:

\[
S(\xi_0,\ \xi_1 \ldots \xi_{n-1}),
\]
then, when we introduce the values (10), we obtain from it a rational function of the radicals

\[
\tau_0,\ \tau_1, \ldots \tau_{n-2}.
\]
But this function does not change when one of these radicals is given another of its $n$ values, and consequently the function must depend rationally on $k_0, k_1 \ldots k_{n-2}$. On account of
%CONT>>
